\documentclass{article}

\usepackage{tabularx}
\usepackage{booktabs}

\title{Reflection and Traceability Report on \progname}

\author{\authname}

\date{}

\input{../Comments}
\input{../Common}

\begin{document}

\maketitle

\plt{Reflection is an important component of getting the full benefits from a
learning experience.  Besides the intrinsic benefits of reflection, this
document will be used to help the TAs grade how well your team responded to
feedback.  Therefore, traceability between Revision 0 and Revision 1 is and
important part of the reflection exercise.  In addition, several CEAB (Canadian
Engineering Accreditation Board) Learning Outcomes (LOs) will be assessed based
on your reflections.}

\section{Changes in Response to Feedback}

The subsections below document \textbf{capstone peer-review} feedback from GitHub issues. For each item we state the source, a concise summary of the issue, our response (what changed), and traceability links. The main pull request bundling MIS, VnV Plan, VnV Report, and related CI workflow updates is \href{https://github.com/hitchly/hitchly/pull/336}{PR \#336}; its description lists \texttt{Closes} for issues \#165, \#167, \#168, \#169, \#244, and \#245 so those issues close when the PR merges to \texttt{main}.

\plt{Summarize the changes made over the course of the project in response to
feedback from TAs, the instructor, teammates, other teams, the project
supervisor (if present), and from user testers.}

\plt{For those teams with an external supervisor, please highlight how the feedback 
from the supervisor shaped your project.  In particular, you should highlight the 
supervisor's response to your Rev 0 demonstration to them.}

\plt{Version control can make the summary relatively easy, if you used issues
and meaningful commits.  If you feedback is in an issue, and you responded in
the issue tracker, you can point to the issue as part of explaining your
changes.  If addressing the issue required changes to code or documentation, you
can point to the specific commit that made the changes.  Although the links are
helpful for the details, you should include a label for each item of feedback so
that the reader has an idea of what each item is about without the need to click
on everything to find out.}

\plt{If you were not organized with your commits, traceability between feedback
and commits will not be feasible to capture after the fact.  You will instead
need to spend time writing down a summary of the changes made in response to
each item of feedback.}

\plt{You should address EVERY item of feedback.  A table or itemized list is
recommended.  You should record every item of feedback, along with the source of
that feedback and the change you made in response to that feedback.  The
response can be a change to your documentation, code, or development process.
The response can also be the reason why no changes were made in response to the
feedback.  To make this information manageable, you will record the feedback and
response separately for each deliverable in the sections that follow.}

\plt{If the feedback is general or incomplete, the TA (or instructor) will not
be able to grade your response to feedback.  In that case your grade on this
document, and likely the Revision 1 versions of the other documents will be
low.} 

\subsection{SRS and Hazard Analysis}

\subsection{Design and Design Documentation}

The following items address \textbf{Module Interface Specification (MIS)} feedback from \textbf{peer reviewers} (GitHub issues). We accepted the feedback and updated \texttt{docs/Design/SoftDetailedDes} (source \texttt{.tex} and rebuilt \texttt{MIS.pdf}); revision history in the MIS documents the March~29 update.

\begin{description}
  \item[\textbf{\#245} --- Notation (\S4) layout and type definitions.]
  \textit{Source:} Peer review (GitHub issue). \\
  \textit{Summary of issue:} The primitive-types table was narrow and hard to read; \texttt{UpdatedFields} used optional-field ``\texttt{?}'' style notation that did not match how partial updates are described elsewhere. \\
  \textit{Response:} We expanded the primitive-types table to full text width using \texttt{tabularx}; we rewrote \texttt{UpdatedFields} in prose so each field may be present or omitted independently, without placeholder syntax. \\
  \textit{Traceability:} Issue \href{https://github.com/hitchly/hitchly/issues/245}{\#245}; changes in \href{https://github.com/hitchly/hitchly/pull/336}{PR \#336} (\path{docs/Design/SoftDetailedDes/sections/04_notation.tex}, \path{docs/Design/SoftDetailedDes/sections/01_revision_history.tex}).

  \item[\textbf{\#244} --- Introduction link to the repository.]
  \textit{Source:} Peer review (GitHub issue). \\
  \textit{Summary of issue:} The MIS should point readers to the implementation repository; the PDF showed only ``can be found at'' because \verb|\href| was missing its visible-text argument. \\
  \textit{Response:} We replaced the link with \verb|\url{https://github.com/hitchly/hitchly}| so the URL is both visible and clickable in the PDF. \\
  \textit{Traceability:} Issue \href{https://github.com/hitchly/hitchly/issues/244}{\#244}; changes in \href{https://github.com/hitchly/hitchly/pull/336}{PR \#336} (\path{docs/Design/SoftDetailedDes/sections/03_introduction.tex}, revision history).
\end{description}

\subsection{VnV Plan and Report}

The following items address \textbf{Verification and Validation Plan} and \textbf{VnV Report} feedback from \textbf{peer reviewers}. We accepted the feedback and updated \texttt{docs/VnVPlan}, \texttt{docs/VnVReport}, and rebuilt the corresponding PDFs; revision-history rows record the March~29 peer-review pass.

\begin{description}
  \item[\textbf{\#169} --- Recording and evaluation of test results.]
  \textit{Source:} Peer review (GitHub issue). \\
  \textit{Summary of issue:} The plan did not clearly state where results are stored, how pass/fail is decided, or how defects are tracked. \\
  \textit{Response:} We added an unnumbered subsection under System Tests (\S3) describing storage (CI, manual logs, metrics), pass/fail rules for automated, manual, and inspection-style tests, and defect handling, with a table-of-contents entry so numbered subsections stay aligned. \\
  \textit{Traceability:} Issue \href{https://github.com/hitchly/hitchly/issues/169}{\#169}; \href{https://github.com/hitchly/hitchly/pull/336}{PR \#336} (\path{docs/VnVPlan/VnVPlan.tex}).

  \item[\textbf{\#165} --- Non-functional system tests (performance, security, scalability).]
  \textit{Source:} Peer review (GitHub issue). \\
  \textit{Summary of issue:} System-level testing for NFRs beyond reliability and usability was thin or missing relative to SRS-style expectations. \\
  \textit{Response:} In the VnV Plan we added \S3.2.3--3.2.5 (performance, security, scalability/capacity) with concrete test cases and tied the \S3 introduction to those areas. In the VnV Report we extended \S4 with matching evaluation areas and test cases (e.g.\ latency, transport/config review, mocked auth/payment, concurrent match load). \\
  \textit{Traceability:} Issue \href{https://github.com/hitchly/hitchly/issues/165}{\#165}; \href{https://github.com/hitchly/hitchly/pull/336}{PR \#336} (\path{docs/VnVPlan/VnVPlan.tex}, \path{docs/VnVReport/sections/04_nonfunctional_requirements.tex}, revision histories).

  \item[\textbf{\#167} --- Traceability in the VnV Plan.]
  \textit{Source:} Peer review (GitHub issue). \\
  \textit{Summary of issue:} The plan's requirements-to-tests mapping needed to cover profile/scheduling (FR-2), ratings/reporting (FR-5), and NFR-3--NFR-5 alongside existing rows. \\
  \textit{Response:} We expanded the \S3.3 traceability table with those requirement IDs, notes for recurring/schedule coverage, and pointers to the new NFR subsections. \\
  \textit{Traceability:} Issue \href{https://github.com/hitchly/hitchly/issues/167}{\#167}; \href{https://github.com/hitchly/hitchly/pull/336}{PR \#336} (\path{docs/VnVPlan/VnVPlan.tex}).

  \item[\textbf{\#168} --- VnV Report trace tables and NFR alignment.]
  \textit{Source:} Peer review (GitHub issue). \\
  \textit{Summary of issue:} Functional vs.\ non-functional captions and matrices did not fully align with SRS G.4 wording; FR-2 recurring/schedule testing and NFR columns were incomplete. \\
  \textit{Response:} We corrected \S9 captions, added a \texttt{test-FR2-3} row for recurring-schedule router coverage, and replaced the NFR table with an NFR-1--NFR-5 matrix consistent with the updated plan. \\
  \textit{Traceability:} Issue \href{https://github.com/hitchly/hitchly/issues/168}{\#168}; \href{https://github.com/hitchly/hitchly/pull/336}{PR \#336} (\path{docs/VnVReport/sections/09_trace_to_requirements.tex}, \path{docs/VnVReport/sections/01_revision_history.tex}).

  \item[\textbf{Repository CI (supporting the same PR).}]
  \textit{Source:} Internal team (merge friction on documentation PRs; not a numbered course issue). \\
  \textit{Summary of issue:} Required GitHub checks stayed on ``Expected --- Waiting for status'' because path-filtered jobs were skipped and, separately, \texttt{cancel-in-progress} on Actions cancelled long LaTeX runs when \texttt{buildtex} auto-pushed PDFs to the PR branch. \\
  \textit{Response:} We updated \path{.github/workflows/ci.yml} so named jobs always report success when paths do not match (no-op steps), fixed the aggregate \texttt{CI} job's result logic, and disabled \texttt{cancel-in-progress} for CI and \texttt{buildtex} to allow runs to finish. Those workflow edits were included on the peer-review PR branch so they reach \texttt{main} when \href{https://github.com/hitchly/hitchly/pull/336}{PR \#336} merges. \\
  \textit{Traceability:} \href{https://github.com/hitchly/hitchly/pull/336}{PR \#336} (\path{.github/workflows/ci.yml}, \path{.github/workflows/buildtex.yml}).
\end{description}

\section{Extras}

\plt{Summarize the extras that were tackled by this project.  Extras
can include usability testing, code walkthroughs, user documentation, formal
proof, GenderMag personas, Design Thinking, etc.  Extras should have already
been approved by the course instructor as included in your problem statement.}

\section{Design Iteration (LO11 (PrototypeIterate))}

For documentation, the Rev$\,\rightarrow\,$Rev1-style iteration on MIS and VnV artifacts followed the course peer-review workflow: feedback was opened as GitHub issues, we grouped related doc edits into a single pull request (\href{https://github.com/hitchly/hitchly/pull/336}{PR \#336}) for coherent review, and we kept revision-history rows and section-level traceability so PDFs remain auditable. We intentionally deferred closing issues until merge (per team workflow) so status stays tied to the PR and \texttt{main}.

\plt{Explain how you arrived at your final design and implementation.  How did
the design evolve from the first version to the final version?} 

\plt{Don't just say what you changed, say why you changed it.  The needs of the
client should be part of the explanation.  For example, if you made changes in
response to usability testing, explain what the testing found and what changes
it led to.}

\section{Design Decisions (LO12)}

\plt{Reflect and justify your design decisions.  How did limitations,
 assumptions, and constraints influence your decisions?  Discuss each of these
 separately.}

\section{Economic Considerations (LO23)}

\plt{Is there a market for your product? What would be involved in marketing your 
product? What is your estimate of the cost to produce a version that you could 
sell?  What would you charge for your product?  How many units would you have to 
sell to make money? If your product isn't something that would be sold, like an 
open source project, how would you go about attracting users?  How many potential 
users currently exist?}

\section{Reflection on Project Management (LO24)}

\plt{This question focuses on processes and tools used for project management.}

\subsection{How Does Your Project Management Compare to Your Development Plan}

\plt{Did you follow your Development plan, with respect to the team meeting plan, 
team communication plan, team member roles and workflow plan.  Did you use the 
technology you planned on using?}

\subsection{What Went Well?}

\plt{What went well for your project management in terms of processes and 
technology?}

\subsection{What Went Wrong?}

\plt{What went wrong in terms of processes and technology?}

\subsection{What Would you Do Differently Next Time?}

\plt{What will you do differently for your next project?}

\section{Ethics and Equity (LO18)}

\plt{Describe how your team applied the principle of equity and universal design
to ensure equitable treatment of all stakeholders.}

\section{Reflection on Capstone}

\plt{This question focuses on what you learned during the course of the capstone project.}

\subsection{Which Courses Were Relevant}

\plt{Which of the courses you have taken were relevant for the capstone project?}

\subsection{Knowledge/Skills Outside of Courses}

\plt{What skills/knowledge did you need to acquire for your capstone project
that was outside of the courses you took?}

\end{document}